\section*{Erd\H{o}s Problem 103: minimum-diameter packings and movable points}

\subsection*{1. Statement and scope}
Let $\Delta(n)$ be the minimum diameter of an $n$-point planar set whose distinct points have mutual distances at least one. Let $h(n)$ be the number of its optimal configurations up to congruence, allowing reflections and ignoring labels. Does $h(n)\to\infty$? This remains unresolved. We prove existence and size bounds, determine the first four orders, and give a sufficient geometric condition for infinitely many optimal configurations at a fixed order.

\subsection*{2. Existence and the diameter scale}
The case $n=1$ has $\Delta(1)=0$. For $n\ge2$, a subset of a $q\times q$ unit grid, with $q=\lceil\sqrt n\rceil$, gives
\[
 \Delta(n)\le\sqrt2(q-1).\tag{1}
\]
Choose a sequence of admissible configurations whose diameters tend to the infimum, label their points, and translate each first point to the origin. Discarding finitely many terms, every point lies in a common bounded closed disk, by (1) and an extra margin of one. A convergent subsequence in the finite product of these disks has a limit configuration. The inequalities $|p_i-p_j|\ge1$ are closed, so the limiting points remain distinct and admissible. Diameter is the maximum of finitely many continuous distance functions, hence is continuous. The limiting set attains $\Delta(n)$.

The open radius-$1/2$ disks about any admissible points are disjoint and lie inside a radius-$(D+1/2)$ disk about one of the points, where $D$ is the diameter. Area comparison gives
\[
 \Delta(n)\ge(\sqrt n-1)/2.\tag{2}
\]
Thus $\Delta(n)=\Theta(\sqrt n)$ as $n\to\infty$. This order-of-magnitude statement does not count optima. Also, every optimum for $n\ge2$ has a pair at distance exactly one: otherwise dilation by the reciprocal of its minimum distance would strictly decrease its positive diameter while preserving admissibility.

\subsection*{3. Four small orders, with the equality case checked}
For $n=1,2,3$ the unique optima up to congruence are respectively one point, a unit segment, and a unit equilateral triangle. Indeed, for $n=2,3$ the diameter is at least one, and attaining one forces every distance to equal one.

For $n=4$, a unit square has diameter $\sqrt2$. We show it is the unique optimum. If the four points form a convex quadrilateral, at least one interior angle is at least $90$ degrees. The two incident sides both have length at least one, so the cosine rule shows that the diagonal joining their other endpoints has length at least $\sqrt2$. If the diameter is exactly $\sqrt2$, no interior angle can exceed $90$ degrees, since then the same inequality is strict. The four angles sum to $360$ degrees, so all are right angles. The rectangle has side lengths at least one and diagonal at most $\sqrt2$, forcing both sides to equal one.

In the remaining cases, if one point lies strictly inside the triangle formed by the other three, the three angular gaps around it are less than $180$ degrees and sum to $360$ degrees. One is at least $120$ degrees. Its two rays end at points of distance at least one from the interior point, and the cosine rule then gives a mutual distance at least $\sqrt3>\sqrt2$. If a point lies on an edge of the triangle, the edge's two parts each have length at least one, so the diameter is at least two. A collinear set of four points has diameter at least three. These cases exhaust the nonconvex configurations. Therefore
\[
 \Delta(4)=\sqrt2,\qquad h(1)=h(2)=h(3)=h(4)=1.\tag{3}
\]

\subsection*{4. A movable-point criterion for infinitely many optima}
Suppose an optimal configuration $P$ of order $n\ge3$ contains a point $v$ satisfying
\[
 1<|v-q|<\Delta(n)\qquad\text{for every }q\in P\setminus\{v\}.\tag{4}
\]
Such a point is sometimes called a rattler; condition (4), rather than terminology, is what is used. Since all its incident distances are strictly below the diameter, a pair of the other, fixed points attains $\Delta(n)$. There are finitely many strict inequalities in (4), so a sufficiently small open disk about $v$ consists entirely of positions satisfying them. Moving $v$ within that disk preserves minimum separation and leaves the fixed diameter pair unchanged. Every resulting configuration is therefore optimal.

These motions yield infinitely many incongruent sets, not just infinitely many coordinate descriptions. Fix two distinct stationary points $a,b\in P\setminus\{v\}$. In any one congruence class of $n$-point sets, the set $S$ of pairwise distance values is finite. If a moved configuration belongs to that class, its moving point $x$ has $|x-a|,|x-b|\in S$. For each pair of radii, the corresponding two circles have at most two intersections. Thus this congruence class accounts for at most $2|S|^2$ possible moving positions. An open disk cannot be covered by finitely many finite sets, so it contains positions in infinitely many distinct congruence classes; in fact there are continuum many.

A useful contrapositive is that if an optimal configuration belongs to an order with finite $h(n)$, every one of its points must be incident to a pair at distance one or to a pair at distance $\Delta(n)$. This is a necessary condition for finiteness, not a proof of rigidity or a sufficient condition for uniqueness.

\subsection*{5. Remaining gap}
No argument above proves that large optimal configurations contain a point satisfying (4), or otherwise produces many optimal shapes. The elementary uniqueness at small orders does not bear on the eventual question. In logical terms, failure of $h(n)\to\infty$ would give a bounded subsequence of values; it would not by itself imply eventual uniqueness or even an eventual uniform bound.

Source: \url{https://www.erdosproblems.com/103}. The existence, packing bounds, equality cases and movable-point criterion are proved in this attempt.

\textbf{Completion rate estimate: 25\%.}

